\section{Planes in Three-Dimensional Space}

A line was produced by allowing one real parameter:
\[
P=P_0+t\mathbf v.
\]
Only one independent direction of motion was available. We now change the
requirement:

\begin{center}
\emph{What geometric set is obtained when motion from one fixed point is allowed
in two independent directions?}
\end{center}

Let $P_0\in\mathbb R^3$ be fixed and let $\mathbf u,\mathbf v$ be nonparallel
vectors. Every permitted displacement is a linear combination
\[
s\mathbf u+t\mathbf v.
\]
Starting at $P_0$ gives
\[
\boxed{
\Pi=\{P_0+s\mathbf u+t\mathbf v:s,t\in\mathbb R\}.
}
\]
One parameter produced a line; two independent parameters produce a plane.
If the two vectors were parallel, both parameters would still move us along only
one direction and no plane would be created.

In coordinates, for
\[
P_0=(x_0,y_0,z_0),\quad
\mathbf u=(u_1,u_2,u_3),\quad
\mathbf v=(v_1,v_2,v_3),
\]
we obtain
\[
\begin{cases}
x=x_0+s u_1+t v_1,\\
y=y_0+s u_2+t v_2,\\
z=z_0+s u_3+t v_3.
\end{cases}
\]

\begin{center}
\begin{tikzpicture}[scale=0.9,>=stealth]
    \coordinate (A) at (0.2,0.8);
    \coordinate (B) at (4.8,0.8);
    \coordinate (C) at (5.9,3.3);
    \coordinate (D) at (1.3,3.3);
    \coordinate (P0) at (2.0,1.7);
    \filldraw[fill=geometricSubsetColor!10,draw=geometricSubsetColor,thick]
        (A)--(B)--(C)--(D)--cycle;
    \draw[blue!70!black,line width=1.1pt,->] (P0)--++(2.0,0)
        node[midway,below] {$\mathbf u$};
    \draw[orange!85!black,line width=1.1pt,->] (P0)--++(0.9,1.1)
        node[midway,left] {$\mathbf v$};
    \fill[geometricSubsetColor] (P0) circle (2.2pt) node[below left] {$P_0$};
    \node[subset label] at (5.0,2.9) {$\Pi$};
\end{tikzpicture}
\end{center}

\subsection{Why a normal vector appears}

The vector chapter proved that
\[
(\mathbf u\times\mathbf v)\cdot(a\mathbf u+b\mathbf v)=0
\]
for every pair of real numbers $a,b$. Therefore
\[
\mathbf n=\mathbf u\times\mathbf v
\]
is perpendicular to every displacement allowed in the plane. This is not an
extra definition imposed on the picture: it is forced by the two generating
directions.

A point $P$ belongs to the plane exactly when the displacement $P-P_0$ is one of
those allowed combinations. Consequently it must satisfy
\[
\boxed{\mathbf n\cdot(P-P_0)=0.}
\]
For $\mathbf n=(a,b,c)$ this becomes
\[
a(x-x_0)+b(y-y_0)+c(z-z_0)=0,
\]
or, after collecting constants,
\[
\boxed{ax+by+cz=d.}
\]
The parametric and normal equations describe the same plane from two different
points of view: one lists allowed motions, while the other states which component
of motion must vanish.

\begin{center}
\begin{tikzpicture}[scale=0.9,>=stealth]
    \coordinate (A) at (0.4,1.0);
    \coordinate (B) at (4.9,1.0);
    \coordinate (C) at (6.0,3.4);
    \coordinate (D) at (1.5,3.4);
    \coordinate (P0) at (2.2,2.0);
    \coordinate (P) at (4.3,2.8);
    \coordinate (N) at (2.2,5.0);
    \filldraw[fill=geometricSubsetColor!10,draw=geometricSubsetColor,thick]
        (A)--(B)--(C)--(D)--cycle;
    \draw[blue!70!black,line width=1.1pt,->] (P0)--(P)
        node[midway,below right] {$P-P_0$};
    \draw[orange!85!black,line width=1.1pt,->] (P0)--(N)
        node[midway,left] {$\mathbf n$};
    \fill[geometricSubsetColor] (P0) circle (2.2pt) node[below left] {$P_0$};
    \fill[geometricSubsetColor] (P) circle (2.2pt) node[above right] {$P$};
    \node[subset label] at (5.2,3.0) {$\Pi$};
\end{tikzpicture}
\end{center}

\subsection{A plane through three points}

Let $P,Q,R$ be three points that do not lie on one line. The displacements
\[
\mathbf u=Q-P,\qquad \mathbf v=R-P
\]
are not parallel. Hence they provide the two independent directions required to
construct a plane:
\[
\boxed{
\Pi=\{P+s(Q-P)+t(R-P):s,t\in\mathbb R\}.
}
\]
Their cross product
\[
\mathbf n=(Q-P)\times(R-P)
\]
is nonzero and perpendicular to both directions, hence to every displacement in
the plane. It therefore supplies the normal equation.

\begin{examplebox}
For
\[
P=(1,0,0),\qquad Q=(0,1,0),\qquad R=(0,0,1),
\]
we obtain
\[
Q-P=(-1,1,0),\qquad R-P=(-1,0,1),
\]
and
\[
(Q-P)\times(R-P)=(1,1,1).
\]
Using the point $P$ gives
\[
(1,1,1)\cdot\bigl((x,y,z)-(1,0,0)\bigr)=0,
\]
so
\[
\boxed{x+y+z=1.}
\]
\end{examplebox}

\subsection{Relative positions of planes}

Let
\[
\Pi_1:\mathbf n_1\cdot P=d_1,
\qquad
\Pi_2:\mathbf n_2\cdot P=d_2.
\]
Their relative position is controlled by the normal directions.

If $\mathbf n_2=\lambda\mathbf n_1$, the planes have the same orientation. They
are either identical or distinct and parallel, depending on the constants. If the
normal vectors are not parallel, the planes meet in a line.

The angle between two planes is measured through their normals. In particular,
\[
\boxed{
\Pi_1\perp\Pi_2
\iff
\mathbf n_1\cdot\mathbf n_2=0.
}
\]

When two nonparallel planes intersect, a direction vector of their intersection
must be perpendicular to both normals. The cross product therefore gives
\[
\boxed{\mathbf v=\mathbf n_1\times\mathbf n_2.}
\]
A point of the intersection is then found by solving the two plane equations.

\subsection{Distance from a point to a plane}

Let $P$ be any point and let $P_0$ be a fixed point on the plane. The shortest
motion from $P$ to the plane must follow the normal direction. Only the component
of $P-P_0$ parallel to $\mathbf n$ contributes to that distance. Hence
\[
\boxed{
d(P,\Pi)=\frac{|\mathbf n\cdot(P-P_0)|}{\|\mathbf n\|}.
}
\]
For $ax+by+cz=d$ and $P=(x_P,y_P,z_P)$ this becomes
\[
\boxed{
d(P,\Pi)=
\frac{|ax_P+by_P+cz_P-d|}{\sqrt{a^2+b^2+c^2}}.
}
\]
The formula is therefore the length of an orthogonal projection, not an isolated
rule.

\begin{summarybox}
A line is generated by one independent direction,
\[
P=P_0+t\mathbf v,
\]
whereas a plane is generated by two,
\[
P=P_0+s\mathbf u+t\mathbf v.
\]
The cross product of the two generating directions produces a normal vector, and
the dot product turns perpendicularity into the equation
\[
\mathbf n\cdot(P-P_0)=0.
\]
Thus the parametric and normal descriptions arise from the same geometric
construction.
\end{summarybox}
